\documentclass[aspectratio=169,11pt]{beamer}

% ==============================================
% PAQUETES
% ==============================================

\usepackage[utf8]{inputenc}
\usepackage[T1]{fontenc}
\usepackage[spanish,es-tabla]{babel}
\usepackage{tgheros}
\usepackage{tikz}

\renewcommand{\familydefault}{\sfdefault}

\hyphenpenalty=10000
\exhyphenpenalty=10000

% ==============================================
% COLORES INSPIRADOS EN REQVIEW
% ==============================================

\definecolor{rvorange}{HTML}{D97728}
\definecolor{rvbrown}{HTML}{734B36}
\definecolor{rvcream}{HTML}{FFF7F0}
\definecolor{rvdark}{HTML}{30343A}
\definecolor{rvmuted}{HTML}{6D737C}
\definecolor{rvline}{HTML}{E9D8CA}

% ==============================================
% ESTILO GENERAL
% ==============================================

\setbeamercolor{normal text}{
  fg=rvdark,
  bg=white
}

\setbeamercolor{frametitle}{
  fg=rvdark,
  bg=white
}

\setbeamercolor{structure}{fg=rvorange}

\setbeamertemplate{navigation symbols}{}

\setbeamersize{
  text margin left=0.65cm,
  text margin right=0.65cm
}

\setbeamerfont{frametitle}{
  size=\Large,
  series=\bfseries
}

% Encabezado
\setbeamertemplate{frametitle}{
  \nointerlineskip
  \begin{beamercolorbox}[
    wd=\paperwidth,
    ht=0.95cm,
    dp=0.15cm,
    leftskip=0.65cm
  ]{frametitle}
    \vfill
    \insertframetitle
    \vfill
  \end{beamercolorbox}

  \hspace*{0.65cm}%
  {\color{rvorange}\rule{1.5cm}{2pt}}
}

% Pie de página
\setbeamertemplate{footline}{
  \leavevmode
  \hbox{
    \begin{beamercolorbox}[
      wd=.82\paperwidth,
      ht=2.7ex,
      dp=1ex,
      leftskip=.65cm
    ]{normal text}
      {\color{rvmuted}\scriptsize
      ISR-401 \quad | \quad GA-SUM-02
      \quad | \quad CES}
    \end{beamercolorbox}

    \begin{beamercolorbox}[
      wd=.18\paperwidth,
      ht=2.7ex,
      dp=1ex,
      rightskip=.65cm plus1fil
    ]{normal text}
      \hfill
      {\color{rvbrown}\scriptsize
      \insertframenumber}
    \end{beamercolorbox}
  }
}

% ==============================================
% TARJETAS FIJAS DE TRES COLUMNAS
%
% Ancho: 3.75 cm
% Alto: 2.65 cm
% ==============================================

\newcommand{\rvcard}[3]{%
\begin{tikzpicture}[baseline=(current bounding box.center)]
  \path[use as bounding box]
    (0,0) rectangle (3.75,2.65);

  \draw[
    draw=rvline,
    fill=rvcream,
    line width=0.7pt,
    rounded corners=5pt
  ] (0.02,0.02) rectangle (3.73,2.63);

  % Número
  \node[
    text=rvorange,
    font=\fontsize{19}{22}\selectfont\bfseries,
    inner sep=0pt
  ] at (1.875,2.17) {#1};

  % Título
  \node[
    text=rvbrown,
    font=\fontsize{11}{13}\selectfont\bfseries,
    text width=3.40cm,
    align=center,
    inner sep=0pt
  ] at (1.875,1.46) {#2};

  % Descripción
  \node[
    text=rvdark,
    font=\fontsize{10}{13}\selectfont,
    text width=3.40cm,
    align=center,
    inner sep=0pt
  ] at (1.875,0.65) {#3};
\end{tikzpicture}%
}

% ==============================================
% TARJETAS FIJAS DE DOS COLUMNAS
%
% Ancho: 5.65 cm
% Alto: 3.50 cm
% ==============================================

\newcommand{\rvwidecard}[4]{%
\begin{tikzpicture}[baseline=(current bounding box.center)]
  \path[use as bounding box]
    (0,0) rectangle (5.65,3.50);

  \draw[
    draw=rvline,
    fill=rvcream,
    line width=0.7pt,
    rounded corners=5pt
  ] (0.02,0.02) rectangle (5.63,3.48);

  % Número
  \node[
    text=rvorange,
    font=\fontsize{18}{21}\selectfont\bfseries,
    inner sep=0pt
  ] at (2.825,2.98) {#1};

  % Título
  \node[
    text=rvbrown,
    font=\fontsize{14}{17}\selectfont\bfseries,
    inner sep=0pt
  ] at (2.825,2.35) {#2};

  % Concepto
  \node[
    text=rvdark,
    font=\fontsize{12}{15}\selectfont\bfseries,
    text width=5.10cm,
    align=center,
    inner sep=0pt
  ] at (2.825,1.57) {#3};

  % Ejemplos
  \node[
    text=rvmuted,
    font=\fontsize{10}{14}\selectfont,
    text width=5.10cm,
    align=center,
    inner sep=0pt
  ] at (2.825,0.65) {#4};
\end{tikzpicture}%
}

% ==============================================
% FILAS HORIZONTALES
%
% \makebox impide que las tarjetas se
% distribuyan en líneas diferentes.
% ==============================================

\newcommand{\rvthree}[9]{%
\noindent
\makebox[\linewidth][c]{%
  \rvcard{#1}{#2}{#3}%
  \hspace{0.20cm}%
  \rvcard{#4}{#5}{#6}%
  \hspace{0.20cm}%
  \rvcard{#7}{#8}{#9}%
}%
}

\newcommand{\rvtwo}[8]{%
\noindent
\makebox[\linewidth][c]{%
  \rvwidecard{#1}{#2}{#3}{#4}%
  \hspace{0.35cm}%
  \rvwidecard{#5}{#6}{#7}{#8}%
}%
}

% ==============================================
% REFERENCIAS BREVES
% ==============================================

\newcommand{\rvsource}[1]{%
  \vfill
  {\fontsize{7}{9}\selectfont
   \color{rvmuted}#1}
}

% ==============================================
% INICIO
% ==============================================

\begin{document}

% ==============================================
% DIAPOSITIVA 1 — PORTADA
% ==============================================

{
\setbeamercolor{background canvas}{bg=rvdark}

\begin{frame}[plain]

\begin{tikzpicture}[remember picture,overlay]

  \fill[rvorange]
    (current page.north west)
    rectangle
    ([yshift=-0.16cm]current page.north east);

  \fill[rvbrown]
    ([xshift=9.6cm]current page.south west)
    -- (current page.south east)
    -- (current page.north east)
    -- cycle;

\end{tikzpicture}

\vspace{0.55cm}

{\color{rvorange}\small\bfseries
UNIVERSIDAD TÉCNICA ESTATAL DE QUEVEDO}

\vspace{0.85cm}

{\color{white}
\fontsize{30}{35}\selectfont
\bfseries Gestión de requisitos}

\vspace{0.30cm}

{\color{rvorange}
\fontsize{26}{30}\selectfont
ReqView}

\vspace{0.50cm}

{\color{white}\large
Sistema de Tutorías y Seguimiento Académico}

\vfill

{\color{white}\small
GA-SUM-02 \quad | \quad Ingeniería de Requisitos}

\vspace{0.30cm}

{\color{white}\fontsize{10}{13}\selectfont
Mayra Córdova \quad · \quad
Angelo Espinoza \quad · \quad
Joseph Saltos}

\vspace{0.45cm}

\end{frame}
}

% ==============================================
% DIAPOSITIVA 2
% ==============================================

\begin{frame}{¿Por qué gestionar requisitos?}

\vspace{0.70cm}

\rvthree
  {01}{DOCUMENTAR}{Registrar necesidades}
  {02}{ORGANIZAR}{Mantener claridad}
  {03}{CONTROLAR}{Gestionar cambios}

\vspace{0.75cm}

\begin{center}
{\large\bfseries\color{rvbrown}
Requisitos claros, desarrollo coherente.}
\end{center}

\rvsource{
Fuente conceptual: IREB,
CPRE Foundation Level, v. 3.3.0 (2026).
}

\end{frame}

% ==============================================
% DIAPOSITIVA 3
% ==============================================

\begin{frame}{Elicitación y gestión}

\vspace{0.48cm}

\rvtwo
  {01}
  {ELICITACIÓN}
  {Descubrir necesidades}
  {Entrevistas, observación y talleres}
  {02}
  {GESTIÓN}
  {Mantener requisitos}
  {Documentación, cambios y trazabilidad}

\vspace{0.60cm}

\begin{center}
{\normalsize\bfseries\color{rvbrown}
ReqView apoya principalmente la gestión.}
\end{center}

\rvsource{
Fuente conceptual: IREB,
CPRE Foundation Level, v. 3.3.0 (2026).
}

\end{frame}

% ==============================================
% DIAPOSITIVA 4
% ==============================================

\begin{frame}{¿Qué es ReqView?}

\vspace{0.25cm}

\begin{center}
{\large\bfseries\color{rvbrown}
Herramienta de gestión de requisitos}
\end{center}

\vspace{0.35cm}

\rvthree
  {01}{ESPECIFICAR}{Documentos estructurados}
  {02}{RELACIONAR}{Trazabilidad}
  {03}{COMPARTIR}{Intercambio ReqIF}

\vspace{0.50cm}

\begin{center}
{\normalsize\color{rvmuted}
Del requisito documentado a su seguimiento}
\end{center}

\rvsource{
Fuente: ReqView,
documentación oficial, www.reqview.com/doc/.
}

\end{frame}

% ==============================================
% DIAPOSITIVA 5
% ==============================================

\begin{frame}{Funcionalidades principales}

\vspace{0.35cm}

\begin{center}
\begin{tikzpicture}

\tikzset{
  feature/.style={
    draw=rvline,
    fill=rvcream,
    rounded corners=5pt,
    minimum width=3.65cm,
    minimum height=1.20cm,
    text width=3.10cm,
    align=center,
    inner sep=0pt,
    font=\fontsize{11}{13}\selectfont\bfseries,
    text=rvbrown
  }
}

\node[feature] at (-2.05,1.45)
  {Documentación};

\node[feature] at (2.05,1.45)
  {Atributos};

\node[feature] at (-2.05,0)
  {Trazabilidad};

\node[feature] at (2.05,0)
  {Filtros};

\node[
  feature,
  fill=rvorange,
  draw=rvorange,
  text=white
] at (0,-1.45)
  {Exportación};

\end{tikzpicture}
\end{center}

\vspace{0.35cm}

\begin{center}
{\normalsize\bfseries\color{rvbrown}
Organizar · Consultar · Relacionar}
\end{center}

\rvsource{
Fuente: ReqView,
documentación oficial, www.reqview.com/doc/.
}

\end{frame}

% Diapositivas 6–9:
% Angelo Espinoza.
 % DIAPOSITIVA 6 — VENTAJAS
 
\begin{frame}{Ventajas para nuestro proyecto}

\vspace{0.70cm}

\rvthree
  {01}
  {CLASIFICAR}
  {Atributos de tipo, prioridad y estado}
  {02}
  {RELACIONAR}
  {Enlaces entre necesidades, requisitos y pruebas}
  {03}
  {REVISAR}
  {Consultar relaciones ante cambios}

\vspace{0.55cm}

\begin{center}
{\normalsize\bfseries\color{rvbrown}
Ejemplo propuesto para el PFC}
\end{center}

\begin{center}
{\small\color{rvdark}
Relacionar el requisito de consulta de indicadores
académicos con su necesidad de origen y su prueba
de acceso autorizado.}
\end{center}

\rvsource{
Funciones: documentación oficial de ReqView,
Custom Requirement Attributes y Requirements Traceability Links.
Criterios de análisis de herramientas:
\cite{ozkaya2023tools}.
}

\end{frame}

 % DIAPOSITIVA 7 — LIMITACIONES Y RIESGOS
 
\begin{frame}{Limitaciones y riesgos}

\vspace{0.48cm}

\rvtwo
  {01}
  {LICENCIA}
  {Comprobar funciones disponibles}
  {La edición gratuita tiene límites.
   Verificar la licencia antes de la demostración.}
  {02}
  {MANTENIMIENTO}
  {La trazabilidad requiere revisión}
  {Los enlaces deben mantenerse.
   La automatización no sustituye el criterio humano.}

\vspace{0.40cm}

\begin{center}
{\normalsize\bfseries\color{rvbrown}
Medida propuesta para el equipo}
\end{center}

\begin{center}
{\small\color{rvdark}
Asignar responsables y revisar las relaciones
cuando cambie un requisito del proyecto.}
\end{center}

\rvsource{
Licencia: ReqView, Frequently Asked Questions.
Esfuerzo y participación humana en trazabilidad:
\cite{ruiz2023traceability}.
}

\end{frame}

% ==============================================
% CONTINUACIÓN DEL EQUIPO
% ==============================================


% Diapositivas 10–14:
% Joseph Saltos.

% Pendiente para la versión grupal:
% - referencias.bib
% - Citas IEEE verificadas
% - Declaración del uso de IA
% - Evidencias originales de ReqView
% - Compilación final reproducible

\end{document}
